% from ega2-1-fr.tex
\oldpage[II]{164}
est, par définition, l'ensemble des idéaux premiers gradués de $\widehat{S}$
contenant $z$, c'est-à-dire $j(X)$ par définition de $j$, ce qui démontre
(\hyperref[II.8.3.2.3-fr]{8.3.2.3}).

Pour prouver enfin la dernière assertion de
(\hyperref[II.8.3.2-fr]{8.3.2}), on peut toujours supposer $Y$ affine. Avec
les notations précédentes, remarquons que dans l'anneau $\widehat{S}$, $z$
n'est pas diviseur de $0$~; comme $\overline{i(C)}=\widehat{C}$, il suffira
de prouver le

\begin{lemma}[8.3.2.7]
\label{II.8.3.2.7-fr}
Soient $T$ un anneau gradué à degrés positifs,
$Z=\operatorname{Proj}(T)$, $g$ un élément homogène de $T$ de degré $d>0$.
Si $g$ n'est pas diviseur de $0$ dans $T$, $Z$ est le plus petit
sous-préschéma fermé de $Z$ qui majore $Z_g=D_+(g)$.
\end{lemma}

En vertu de (\hyperref[I.4.1.9-fr]{I, 4.1.9}) la question est locale sur $Z$~;
pour tout élément homogène $h\in T_e$ ($e>0$), il suffit donc de prouver que
$Z_h$ est le plus petit sous-préschéma fermé de $Z_h$ qui majore $Z_{gh}$~;
il résulte des définitions et de (\hyperref[I.4.3.2-fr]{I, 4.3.2}) que cette
condition équivaut au fait que l'homomorphisme canonique
$T_{(h)}\to T_{(gh)}$ est \emph{injectif}. Or cet homomorphisme s'identifie à
l'homomorphisme canonique
$T_{(h)}\to(T_{(h)})_{g^e/h^d}$ (\hyperref[II.2.2.3-fr]{2.2.3}). Mais comme
$g^e$ n'est pas diviseur de $0$ dans $T$, $g^e/h^d$ ne l'est pas dans $T_h$
(ni \emph{a fortiori} dans $T_{(h)}$), car la relation
$(g^e/h^d)(t/h^m)=0$ avec $t\in T$ et $m>0$ entraîne l'existence d'un $n>0$
tel que $h^ng^et=0$, d'où $h^nt=0$, et par suite $t/h^m=0$ dans $T_h$. Cela
achève donc la démonstration (\hyperref[0.1.2.2-fr]{0, 1.2.2}).

\begin{env}[8.3.3]
\label{II.8.3.3-fr}
On identifiera souvent le cône affine $C$ au sous-préschéma induit par le cône
projectif $\widehat{C}$ sur l'ouvert $i(C)$, au moyen de l'immersion ouverte
$i$. Le sous-préschéma fermé de $C$, associé à l'immersion fermée
$\varepsilon$, est appelé le \emph{préschéma des sommets de $C$}~; on dit
aussi que $\varepsilon$, qui est une $Y$-section de $C$, est la \emph{section
sommet} ou la \emph{section nulle de $C$}~; on peut identifier $Y$ au
préschéma des sommets de $C$ au moyen de $\varepsilon$. D'ailleurs
$i\circ\varepsilon$ est une $Y$-section de $\widehat{C}$, donc aussi une
immersion fermée (\hyperref[I.5.4.6-fr]{I, 5.4.6}), correspondant à
l'homomorphisme canonique surjectif de degré $0$,
$\widehat{\mathcal{S}}=\mathcal{S}[z]\to\mathcal{O}_Y[z]$
(cf. (\hyperref[II.3.1.7-fr]{3.1.7})), de noyau
$\mathcal{S}_+[z]=\mathcal{S}_+\widehat{\mathcal{S}}$~; le sous-préschéma de
$\widehat{C}$ associé à cette immersion fermée est encore appelé le
\emph{préschéma des sommets de $\widehat{C}$}, et $i\circ\varepsilon$ la
\emph{section sommet de $\widehat{C}$}~; il peut s'identifier par
$i\circ\varepsilon$ à $Y$. Enfin, le sous-préschéma fermé de $\widehat{C}$
associé à $j$ est appelé le \emph{lieu à l'infini de $\widehat{C}$}, et peut
être identifié à $X$ au moyen de $j$.
\end{env}

\begin{env}[8.3.4]
\label{II.8.3.4-fr}
Les sous-préschémas de $C$ (resp. $\widehat{C}$) induits respectivement sur
les \emph{ouverts}
\begin{equation}
\label{II.8.3.4.1-fr}
  E=C-\varepsilon(Y),\qquad
  \widehat{E}=\widehat{C}-i(\varepsilon(Y))
\tag{8.3.4.1}
\end{equation}
sont appelés respectivement (par abus de langage) le \emph{cône affine
épointé} et le \emph{cône projectif épointé} définis par $\mathcal{S}$~; on
notera qu'en dépit de cette terminologie, $E$ \emph{n'est pas nécessairement
affine} sur $Y$, ni $\widehat{E}$ projectif sur $Y$
(cf. (\hyperref[II.8.4.3-fr]{8.4.3})). Lorsqu'on identifie $C$ à $i(C)$, on a
donc pour les espaces sous-jacents
\begin{equation}
\label{II.8.3.4.2-fr}
  C\cup\widehat{E}=\widehat{C},\qquad C\cap\widehat{E}=E
\tag{8.3.4.2}
\end{equation}
de sorte que $\widehat{C}$ peut être considéré comme obtenu par
\emph{recollement} des sous-préschémas ouverts $C$ et $\widehat{E}$~; en
outre, en vertu de (\hyperref[II.8.3.2.3-fr]{8.3.2.3}),
\begin{equation}
\label{II.8.3.4.3-fr}
  E=\widehat{E}-j(X).
\tag{8.3.4.3}
\end{equation}

